% Published transcription: DKY22, print pp.402--404, physical pp.40--42.
% Only the energy-pairing definition and full Section 9.2 are retained.
% Original numbered references are retained literally: [19]=FRL06; [20]=Fil17.
\section*{Published author argument reused in Fu's Theorem 3.13}
\noindent Laura DeMarco, Holly Krieger and Hexi Ye, \emph{Common preperiodic points for quadratic polynomials}, J. Mod. Dyn. 18 (2022), 363--413, \doi{10.3934/jmd.2022012}, pp.402--404. Original reference [19] is Favre--Rivera-Letelier [FRL06], and [20] is Fili [Fil17] in the retained bibliography. This appendix transcribes their published author text; Fu's own substitutions and proof remain verbatim in Section 3 above.

There is a pairing between any two such heights, $h_{\mu}$ and $h_{\nu}$, associated to adelic measures $\mu$ and $\nu$, as
\begin{equation}
\langle h_{\mu},h_{\nu}\rangle=\frac12\sum_{v\in M_K}\frac{[K_v:\Q_v]}{[K:\Q]}(\mu_v-\nu_v,\mu_v-\nu_v)_v.\tag{9.2}\label{DKY22:eq9.2}
\end{equation}
It satisfies $\langle h_{\mu},h_{\nu}\rangle=0\Longleftrightarrow h_{\mu}=h_{\nu}\Longleftrightarrow\mu=\nu$. The energy pairing (1.3) between two quadratic polynomials is a special case, taking the dynamical canonical heights $\hat h_{c_1}$ and $\hat h_{c_2}$ associated to their adelic equilibrium measures.

\subsection*{9.2. Height pairing as a distance}
Following [20], we consider a distance between two adelic measures $\mu=\{\mu_v\}$ and $\nu=\{\nu_v\}$ on $\PP^1$ over a number field $K$, defined by
\[
d(\mu,\nu):=\langle h_{\mu},h_{\nu}\rangle^{1/2},
\]
where $\langle h_{\mu},h_{\nu}\rangle$ was defined in (9.2); see [20, Theorem 1].

Suppose that $c_1$ and $c_2$ are elements of a number field $K$. Let $\mu_1:=\{\mu_{c_1,v}\}_{v\in M_K}$ and $\mu_2:=\{\mu_{c_2,v}\}_{v\in M_K}$ be the equilibrium measures of $f_{c_1}$ and $f_{c_2}$, respectively. Let $F$ be any finite, nonempty, $\Gal(\bar K/K)$-invariant subset of $\PP^1(\bar\Q)$. Let $[F]$ denote the probability measure supported equally on the elements of $F$. For each place $v$ of $K$, choose a positive real $\varepsilon_v>0$, with $\varepsilon_v=1$ for all but finitely many $v$. The collection $\varepsilon:=\{\varepsilon_v\}_{v\in M_K}$ will be called an adelic radius. As in [19], we consider the adelic measure $[F]_{\varepsilon}$, defined as a regularization of the probability measure $[F]$: it is supported on the circles of radius $\varepsilon_v$ about each point of $F$. At a non-archimedean place, this means the Type II or III point associated to the disk of radius $\varepsilon_v$. The triangle inequality implies that
\begin{equation}
\langle f_{c_1},f_{c_2}\rangle^{1/2}=d(\mu_1,\mu_2)\leq d(\mu_1,[F]_{\varepsilon})+d(\mu_2,[F]_{\varepsilon})\tag{9.3}\label{DKY22:eq9.3}
\end{equation}
for any choices of $F$ and $\varepsilon$.

It is worth noting that, if the radius $\varepsilon_v\to0$ at some place, then the right-hand-side of (9.3) will tend to $\infty$. This is because the potential of the measure $[F]_{\varepsilon}$ at $v$ will blow up near the points of $F$. On the other hand, for $\varepsilon_v$ too large, the measure $[F]_{\varepsilon}$ is not a good approximation of $[F]$. Thus, for (9.3) to be useful in our proof of Theorem 1.1, we will need to choose $\varepsilon$ well. This general strategy also appears in the proofs of [19, Theorem 3] and in [20, Proposition 13]. In our case, the choice of $\varepsilon=\{\varepsilon_v\}_{v\in M_K}$ will be governed by Proposition 3.8 and its non-archimedean counterparts, and this leads to Theorem 1.8.

\begin{dkylemma}
Let $K$ be a number field and fix $c_1\neq c_2$ in $K$. We have
\[
\langle f_{c_1},f_{c_2}\rangle^{1/2}\leq\sum_{i=1}^{2}\left(\sum_{v\in M_K}\frac{[K_v:\Q_v]}{[K:\Q]}\left(-(\mu_i,[F]_{\varepsilon})_v+\frac{\log(1/\varepsilon_v)}{2|F|}\right)\right)^{1/2}
\]
for any choice of finite, non-empty, $\Gal(\bar K/K)$-invariant subset $F$ of $\bar\Q$ and any adelic radius $\varepsilon=\{\varepsilon_v\}_{v\in M_K}$.
\end{dkylemma}

\begin{proof}
We first observe that
\begin{align*}
d(\mu_i,[F]_{\varepsilon})^2&=\frac12\sum_v\frac{[K_v:\Q_v]}{[K:\Q]}(\mu_i-[F]_{\varepsilon},\mu_i-[F]_{\varepsilon})_v\\
&=\sum_v\frac{[K_v:\Q_v]}{[K:\Q]}\left(-(\mu_i,[F]_{\varepsilon})_v+\frac12([F]_{\varepsilon},[F]_{\varepsilon})_v\right)
\end{align*}
because $(\mu_i,\mu_i)_v=0$ at every place. The self-pairing of $[F]_{\varepsilon}$ can be estimated in terms of the self-pairing of $[F]$ ([20, Lemma 12] and [19, Lemma 4.11]), as
\[
([F]_{\varepsilon},[F]_{\varepsilon})_v\leq([F],[F])_v+\frac{\log(1/\varepsilon_v)}{|F|}.
\]
But observe that
\[
\sum_v\frac{[K_v:\Q_v]}{[K:\Q]}([F],[F])_v=0
\]
by the product formula on $K$. So the triangle inequality (9.3) completes the proof of the proposition.
\end{proof}
